\subsection{Résultats}
\label{sec:DSCresultats}


Lors de la chauffe rapide (Fig. \ref{fig:chauffe_rapide}), nous constatons une température de transition vitreuse autour de 82 °C et un pic exothermique vers 244,8 °C d’une intensité de -43,15 J/g. En effet, chauffé à froid, les cristaux qui se défont libèrent une certaine chaleur. En revanche, le refroidissement rapide à chaud sur la Fig. \ref{fig:refroidissement_rapide} ne laisse pas au matériau le temps de se restructurer correctement et prend une structure amorphe.



\begin{flushleft}
\begin{minipage}{0.49\textwidth} 
\includegraphics[width=\textwidth]{dsc_chauffe.pdf}%
	\captionof{figure}{$\delta$H(T[$^\circ$C]) chauffe rapide}
       \label{fig:chauffe_rapide}
\end{minipage}
\hfill
\begin{minipage}{0.49\textwidth}
    \includegraphics[width=\linewidth]{dsc_amorphe.pdf}
	\captionof{figure}{$\delta$H(T[$^\circ$C]) refroidissement rapide}
        \label{fig:refroidissement_rapide}
\end{minipage}
\end{flushleft}
\\\\
Par la suite, le comportement de l’échantillon change lorsque nous lui permettons de se refroidir ou de chauffer plus lentement (visible sur les Figs. \ref{fig:chauffe_lente} et \ref{fig:refroidissement_lent}).

\begin{flushleft}
\begin{minipage}{0.49\textwidth}
    \includegraphics[width=\linewidth]{dsc_cristallisation_chaude.pdf}
        \captionof{figure}{$\delta$H(T[$^\circ$C]) chauffe lente}
        \label{fig:chauffe_lente}
\end{minipage}
\hfill
\begin{minipage}{0.49\textwidth}
    \includegraphics[width=\linewidth]{dsc_cristallisation_froide.pdf}
        \captionof{figure}{$\delta$H(T[$^\circ$C]) refroidissement lent}
        \label{fig:refroidissement_lent}
\end{minipage}
\end{flushleft}
Après un refroidissement rapide sur la Fig. \ref{fig:refroidissement_rapide}, nous chauffons lentement et obtenons la Fig. \ref{fig:chauffe_lente} sur laquelle $T_v$ se situe autour de 77 °C, soit environ 5 °C plus bas que lors de la première chauffe sur la Fig. \ref{fig:chauffe_rapide}. Nous observons ensuite un pic endothermique sur la Fig. \ref{fig:chauffe_lente} autour de 131 °C que nous ne voyons pas sur la Fig. \ref{fig:chauffe_rapide}. Cela démontre l’énergie de recristallisation du matériau chauffé à froid suite au refroidissement rapide ayant laissé une structure amorphe. Autour de 245 °C, avec 1,1 °C de différence par rapport à la chauffe rapide (Fig. \ref{fig:chauffe_rapide}), nous observons un second pic sur la Fig. \ref{fig:chauffe_lente}, qui est également de nature exothermique. Cela démontre l’énergie libérée lors de la décristallisation dans l’échantillon de PET. Lors du refroidissement lent sur la Fig. \ref{fig:refroidissement_lent}, nous observons, contrairement à la Fig. \ref{fig:refroidissement_rapide}, un pic endothermique démontrant l’énergie nécessaire à la recristallisation du matériau.

\subsection{Analyse des résultats}

Dans son ensemble, chaque expérience en DSC est visible dans le Tab. \ref{val:peaks}, dont on connaît alors l’enthalpie de fusion à partir de la mesure de l’aire.

\begin{table}[!ht]
\centering
\footnotesize
\setlength{\tabcolsep}{4pt}
\renewcommand{\arraystretch}{0.9}
\begin{tabular}{|c|c|c|c|c|c|}
\hline
        \textbf{Figure}
        & \textbf{$V_v_t$ [$K/min$]}
        & \textbf{Pic complexe [$^\circ$]}
        & \textbf{Aire [$J/g$]}
        & \textbf{$T_v$ [bool]}
	& \textbf{$T_f$ [$^\circ$C]}
\\\hline
	\ref{fig:chauffe_rapide} & 50.0 & 244.8 & -43.15 & Oui & 68.4 \\
	\ref{fig:refroidissement_rapide} & -200.0 & $\emptyset$ & $\emptyset$ & $\emptyset$ & $\emptyset$\\
	\ref{fig:chauffe_lente} & 10.0 & 131.1-245.9 & 23.93-(-35.58) & Oui & 77.7 \\
	\ref{fig:refroidissement_lent} & -10.0 & 186.8 & 36.19 & $\emptyset$ & $\emptyset$ \\
\hline
\end{tabular}
        \caption{$V_v_t$: vitesse de la variation thermique, $T_v$: Transition vitreuse}
\label{val:peaks}
\end{table}
La valeur théorique d’enthalpie de fusion du PET est définie par la Ref. \cite{petnet} à 140 J/g. Pour la chauffe lente (chauffe à froid) sur la Fig. \ref{fig:chauffe_lente}, nous obtenons la valeur mesurée d’enthalpie de fusion à -43,15 J/g à 244,8 °C. Ainsi, nous obtenons le taux de cristallinité en fonction de la rampe thermique de 50.0 $K/min$ $X_2_4_4_._8$(50.0) :


\[
	X_2_4_4_._8(50.0) = |\frac{\delta H_f_t_h}{\delta H_f_m_e}| \cdot 100\% = \frac{35,58}{140} \cdot 100 \% = 25.41 \%.
\]
Pour la même température, à ~1 °C près, nous obtenons pour la chauffe rapide sur la Fig. \ref{fig:chauffe_rapide} le taux de cristallinité $X_2_4_5_._9$(10.0) :

\[
	X_2_4_5_._9(10.0) = |\frac{\delta H_f_t_h}{\delta H_f_m_e}| \cdot 100\% = \frac{43.15}{140} \cdot 100\% = 32.25 \%.
\]

Nous constatons que $X_2_4_5_._9$(10.0) $>$ $X_2_4_5_._9$(50.0), c’est-à-dire que le taux de cristallinité avant décristallisation est plus élevé dans l’échantillon chauffé lentement que celui chauffé rapidement. En effet, suite à avoir été chauffé rapidement, nous observons une phase amorphe lors du refroidissement rapide sur la Fig. \ref{fig:refroidissement_rapide}. Cela veut dire que le PET n’a pas eu le temps de récupérer le premier taux de cristallinité mesuré avec la chauffe rapide. Nous pouvons vérifier quel était le taux de cristallinité à 131,1 °C observé sur la chauffe lente sur la Fig. \ref{fig:chauffe_lente} suite au refroidissement amorphe de l’échantillon avec $X_1_3_1_._1$(10.0) :

\[
	X_1_3_1_._1(10.0) = |\frac{\delta H_f_t_h}{\delta H_f_m_e}| \cdot 100\% = \frac{23.93}{140} \cdot 100\% = 17.11 \%,
\]

Nous remarquons que le taux de cristallisation est inférieur au taux de décristallisation, soit $X_1_3_1_._1$(10.0) $<$ $X_2_4_5_._9$(10.0). Le matériau a, en fin de compte, un taux de cristallinité plus élevé quand il est chauffé, puis refroidi rapidement avant d’être rechauffé lentement de 50 °C à 300 °C.

\\\\
Enfin, quand le matériau est à nouveau refroidi, mais plus lentement voir Tab. \ref{val:peaks}, nous obtenons $X_1_8_6_._8$(-10.0), soit : 

\[
        X_1_8_6_._8(-10.0) = |\frac{\delta H_f_t_h}{\delta H_f_m_e}| \cdot 100\% = \frac{36.19}{140} \cdot 100\% = 25.85 \%.
\]

Nous obtenons par conséquent pratiquement le même taux de cristallinité avec $X_1_8_6_._8$(-10.0) que pour $X_2_4_5_._9$(10.0), soit par une différence de 0.44 \%.

\section{Conclusion}
Pour conclure, nous avons pu déterminer grâce aux essais de DMA l’énergie d’activation Ea
de 295.730 ± 18.959[kJ/mol] en se basant sur la dépendance des modules viscoélastiques et
de la température de transition vitreuse avec la fréquence d’excitation d’une masse oscillante
sur notre échantillon de PMMA.

\\\\
Nous avons pu identifier sur notre échantillon de PET et grâce aux essais de DSC ses diffé-
rentes transitions thermiques en fonction de plusieurs rampes thermiques. Les températures
de transitions vitreuses obtenues furent après un refroidissement rapide : Tv (10) = 77$^\circ$C
et chauffé initialement et rapidement à froid : Tv (50) = 82$^\circ$C. Les températures de fusion
furent : Tf (50) = 68.4 $^\circ$C puis Tf (10) = 77,7 $^\circ$. Enfin, grâce aux aires mesurées nous pûmes iden-
tifier différents taux de cristallinité de l’échantillon. Ces derniers allant de 17.11 \% à 32.25 \%.
\\\\
Nous pouvons envisager de procéder à des mesures avec une plage de fréquence plus large
pour obtenir une énergie d’activation plus précise du PMMA. Pour le DSC, nous pou-
vons imaginer de faire varier les rampes thermiques et tester plusieurs autres combinaisons
chauffe/refroidissement afin d’en apprendre plus sur les propriétés thermiques de l’échantillon
de PET et comment le taux de cristallinité de l’échantillon varie en fonction de plusieurs en-
vironnements thermiques consécutifs.
\\\\
Enfin, nous avons appris que les techniques de DMA et DSC sont complémentaires afin de
pouvoir caractériser les propriétés thermiques d’un polymère. Que ce soit pour un échantillon
de PMMA, PET, ABS, TPU, PEEK et tant d’autres parmi cette concaténation organique
des mailles des polymères.

